% chapter/chapter3.tex
% The Cursed & Corrupted – Lycanthropes, Vampires, and the Children of Malachai
% Rewritten in the voice of Lerris Fenwood, with marginalia from the Grey Wanderer,
% Saikou Ira, and Dusana of the Raven Road.

\epigraph{“Malachai does not curse. She offers gifts that remember to be hungry. The coin that never spends is the one you don’t remember taking.”}{—The Grey Wanderer, from a conversation in a Zakov tavern}

The curse is not a punishment. I learned this in the catacombs beneath Zahk’Kozahn, where a woman who had been a werewolf for forty years asked me for a glass of water. She did not drink it. She held it to her forehead and wept. The water turned to salt. The salt was warm.

Malachai’s gifts are not evil. They are **desperate**. They are the answer you receive when you ask for help from the wrong threshold. This chapter catalogues those who have accepted a gift they cannot return: lycanthropes, vampires, slashers, memory‑eaters, and the hollowed things that wander the cursed city. I have met some of them. I have pitied all of them.

\section*{Lycanthropes – The Shape of Hunger}

\noindent Lycanthropy is a spiritual infection – a fragment of Malachai's chained essence passed from host to host through violence, but only if the victim is in a state of profound emotional vulnerability. It is not a disease; it is a **story** that the body learns to tell.

\subsection*{Linn Sea‑Wolf}

\noindent\textbf{Description}: Transforms into a massive seal or walrus. Solitary, melancholic, preys on ships. The transformation is triggered by the first winter storm after a kill.

\noindent\textbf{Stats}: Tier III, Cap 4. Shapechanger, Amphibious, Solitary.\\
\textbf{Harm}: 5. \textbf{Fatigue}: Body 4.

\noindent\textbf{Traits}: 
\emph{Claw} (DV 2, Harm 2), \emph{Drowning Grasp} (DV 4, Harm 1 + Suffocation). \emph{Silver Vulnerability} (+2 Harm from silver). \emph{Transformation Trigger} – first winter storm after a kill. The storm does not cause the change; it **permits** it.

\noindent\textbf{Resolution}: Silver weapon to the heart or brain. Aconite poultice can prevent transformation if applied within the hour. The Well of Ichor can purge lycanthropy, but the journey to Zahk’Kozahn is itself a trial.

\noindent\textbf{Clock}: \emph{Moon’s Hunger [6]} – when filled, the lycanthrope enters a frenzy. The frenzy does not end until something is dead. The Sea‑Wolf prefers to drown.

\saikou{I tracked a Sea‑Wolf for three weeks. She had been a fisher’s wife. When I finally found her, she was sitting on a rock, watching the waves. She said, “I remember my daughter’s face. That is the curse. Not the teeth.” I left her there.}

\subsection*{Ykrul Sky‑Hound}

\noindent\textbf{Description}: Hawk‑headed hunter that transforms during the day. Hunts oath‑breakers, not flesh. The transformation is triggered by the sight of blood on fresh snow.

\noindent\textbf{Stats}: Tier III, Cap 4. Shapechanger, Aerial, Solitary.\\
\textbf{Harm}: 4. \textbf{Fatigue}: Body 3.

\noindent\textbf{Traits}: 
\emph{Claw} (DV 2, Harm 1), \emph{Blood Scent} (automatically detects oath‑breakers within Far). \emph{Silver Vulnerability} (+2 Harm from silver). \emph{Transformation Trigger} – sight of blood on fresh snow. The snow does not need to be clean; the blood does not need to be human.

\noindent\textbf{Resolution}: Silver weapon, or a public confession of the oath broken. The Sky‑Hound will accept the confession as payment. I have seen a man confess to a lie he told twenty years ago; the Hound nodded and flew away.

\noindent\textbf{Clock}: \emph{Blood Debt [6]} – ticks each time the Hound hunts without finding an oath‑breaker. When full, it turns on its own pack.

\grey{The Ykrul do not fear the Sky‑Hound. They fear being the reason it hunts.}

\subsection*{Valewood Root‑Walker}

\noindent\textbf{Description}: Transforms into a tree – rooted, immobile, but all‑seeing. Guardians, not predators. The transformation is triggered by a witnessed promise of hospitality.

\noindent\textbf{Stats}: Tier III, Cap 4. Shapechanger, Immobile, All‑Seeing.\\
\textbf{Harm}: 6 (but cannot move to attack). \textbf{Fatigue}: Body 5.

\noindent\textbf{Traits}: 
\emph{Root Grasp} (DV 3, Harm 1 + Entangle), \emph{Forest Sight} (sees everything within Far range of its location). \emph{Silver Vulnerability} (+1 Harm from silver). \emph{Transformation Trigger} – a witnessed promise of hospitality. The promise must be broken for the transformation to complete.

\noindent\textbf{Resolution}: Fulfill the original promise of hospitality. The Root‑Walker will return to human form and thank you. If the promise was made to a dead person, the Root‑Walker will accept a substitution: a gift to the dead’s family, a story told at their grave.

\noindent\textbf{Clock}: \emph{Root’s Patience [8]} – ticks each time the broken promise is ignored. When full, the Root‑Walker becomes a permanent tree. It will not die. It will only wait.

\dusana{The Root‑Walker in the Valewood was a woman who promised to shelter a stranger. The stranger died on her doorstep. She planted an acorn in his hand and became the tree. The acorn grew. The woman still watches.}

\subsection*{Aeler Stone‑Hide}

\noindent\textbf{Description}: Massive, slow, nearly invulnerable. Transforms only underground. A mining hazard. The transformation is triggered by a cave‑in.

\noindent\textbf{Stats}: Tier III, Cap 5. Shapechanger, Sturdy, Subterranean.\\
\textbf{Harm}: 7. \textbf{Fatigue}: Body 5.

\noindent\textbf{Traits}: 
\emph{Claw} (DV 2, Harm 2), \emph{Crush} (DV 4, Harm 3). \emph{Silver Vulnerability} (+1 Harm from silver; silver is scarce underground). \emph{Transformation Trigger} – a cave‑in. The Stone‑Hide does not cause the cave‑in; it is caused **by** the cave‑in. It is the mountain’s immune response.

\noindent\textbf{Resolution}: Silver weapon, or a **keystone clause**. If the original cave‑in was caused by a mining violation, correcting the violation (shoring the tunnel, paying wergild to the dead) will reverse the transformation. The Stone‑Hide will weep stone dust and sleep.

\noindent\textbf{Clock}: \emph{Stone’s Grudge [6]} – ticks each time the transformation is reversed without addressing the original cave‑in. When full, the Stone‑Hide becomes a permanent part of the mountain. It will not leave; it will not forget.

\lerris{I have seen a Stone‑Hide weep. The tears were not water. They were fine, grey dust. The dust settled on a ledger and erased a name. The name was the miner who had failed to count his breaths.}

\section*{Vampires – The First Children}

\noindent Vampires are the eldest children of Malachai, the first to drink from the Chained Angel’s wound. The Thirteen First Vampires still walk the world, each bound to return to Zahk’Kozahn once per century to pay tribute. Their spawn are weaker, but no less hungry.

\subsection*{The Thirteen (Abridged)}

I have met three of the Thirteen. I will describe them briefly, as a warning.

\noindent\textbf{Baron Quentin Asselin}: The first to drink. Disillusioned, melancholic, almost gentle. Converts guests or drives them mad. Never accept his wine.

\noindent\textbf{Sister Maris}: Runs a children’s hospice. The children never leave. The Everflame has tried to destroy her hospice four times. Each time, the children defended her.

\noindent\textbf{The Sparrow}: Kills only the wicked, or so he claims. Cannot look at mirrors. Has saved my life twice. I suspect he is merely bored.

\saikou{The Sparrow once left a silver coin on my desk. The coin had a scratch that looked like a name. I did not spend it. I kept it in a box. The box is now warm.}

\subsection*{Lesser Bloodlines}

\noindent\textbf{Linn Draugr}: Breathe underwater, trapped by sea‑salt circles. Weakness: a circle of salt poured at high tide will hold until dawn.

\noindent\textbf{Aeler Stone‑Sleepers}: Drink breath, awakened by bronze bells at dawn. Weakness: a bell rung at the entrance to their lair will paralyse them for one exchange.

\noindent\textbf{Ykrul Bone‑Eaters}: Turn to dust when killed, but rise if their name is spoken. Erase the name. I have seen a man kill a Bone‑Eater seven times. The eighth time, he burned the body and scattered the ash in a river. The Bone‑Eater did not return.

\noindent\textbf{Vhasian Candle‑Wights}: Extinguish flames through windows. A lit candle at the threshold bars them. Weakness: a candle lit in a silver holder will burn for three hours and keep them at bay.

\noindent\textbf{Lethai‑ar Silk‑Drinkers}: Feed on desire. Vulnerable to vows witnessed by a third party. Weakness: a vow spoken aloud in the presence of a neutral witness will bind the Silk‑Drinker for a night.

\subsection*{Vampire (Lesser Bloodline) – Template}

\noindent\textbf{Stats}: Tier III, Cap 4. Undead, Hypnotic Gaze, Blood Drain.\\
\textbf{Harm}: 5. \textbf{Fatigue}: Body 4.

\noindent\textbf{Traits}: 
\emph{Claw} (DV 2, Harm 1), \emph{Blood Drain} (DV 3, Harm 1 + 1 Fatigue to victim, heals vampire 1 Harm). \emph{Sunlight Vulnerability} – 1 Harm per round in direct sun. \emph{Invitation Required} – cannot enter a home without an invitation.

\noindent\textbf{Resolution}: Stake through the heart (DV 5, must be wooden), decapitation, fire. A silver mirror held to the face will force a Resolve test (DV 4) or the vampire will flee.

\noindent\textbf{Clock}: \emph{Thirst [6]} – when filled, the vampire loses control and feeds openly. The feeding is always witnessed. The witness is always the next victim.

\grey{Do not pity the vampire. Pity the hunger. The vampire is only its symptom.}

\section*{Other Children of Malachai – The Cursed}

\noindent Not all of Malachai’s children are vampires or werewolves. Her gifts manifest according to the desires and sins of those who accept them.

\subsection*{The Unkillable Killer (Slasher)}

\noindent\textbf{Description}: A mortal who bargained for invincibility – to survive a massacre, to avenge a family, to become the hunter instead of the prey. The curse: they cannot stop killing. The hunger for violence grows with each life taken.

\noindent\textbf{Stats}: Tier II‑III, Cap 4. Corporeal, Regenerating, Hungry.\\
\textbf{Harm}: 6. \textbf{Fatigue}: Body 4.

\noindent\textbf{Traits}: 
\emph{Regeneration} – ignore first 2 Harm each scene. \emph{Unkillable} – cannot die from mundane harm. Decapitation, fire, and magical destruction may kill (GM discretion). \emph{The Face in the Mirror} – once per session, GM may force a Resolve test (DV 4) or suffer −1 die to all rolls for the scene.

\noindent\textbf{Resolution}: Silver, fire, or finding the grave of the first victim. The face in the mirror is the first victim. If the killer speaks their name, the regeneration stops for one exchange.

\noindent\textbf{Clock}: \emph{Hunger Clock [6]} – ticks each time the killer kills a sentient being. At 3, social rolls −1 die. At 6, frenzy – attack nearest creature.

\lerris{I met a Slasher in the sewers of Zakov. He had been a dockworker. He asked me to read a letter he had written to his daughter. I read it. He wept. Then he killed a rat and ate it. The hunger is not a choice; it is a tax.}

\subsection*{The Memory Eater (Psychic Vampire)}

\noindent\textbf{Description}: A scholar who wanted to know everything; a lover who wanted to never forget; a spy who wanted to steal secrets without trace. The curse: they must feed on fresh memories or lose their own.

\noindent\textbf{Stats}: Tier II, Cap 3. Corporeal, Telepathic, Hungry.\\
\textbf{Harm}: 3. \textbf{Fatigue}: Spirit.

\noindent\textbf{Traits}: 
\emph{Memory Feed} (touch a willing or helpless target, take one specific memory; target forgets permanently; gain 1 Boon). \emph{Emotion Vampire} (+1 die after feeding). \emph{Phasing} (intangible 1 round per scene).

\noindent\textbf{Resolution}: Feeding on a false memory disorients it. The Memory Eater cannot distinguish a true memory from a well‑constructed lie. A skilled deceiver can feed it a false history; the Eater will consume it and become confused, suffering −1 die to all rolls for a scene.

\noindent\textbf{Clock}: \emph{Hunger for Emotion [6]} – must feed on a significant memory once per week or lose 1 Fatigue per day. The fatigue is not physical; it is the slow erosion of identity.

\saikou{I once fed a Memory Eater a false memory of a murder that never happened. The Eater spent three days trying to confess. The confession was beautiful, detailed, and entirely invented. When it realised the truth, it wept and vanished. I have not seen it since.}

\subsection*{The Plague Bearer (Contagion)}

\noindent\textbf{Description}: A healer who wanted to cure a plague; a mother who wanted to spare her child; a martyr who wanted to take suffering onto themselves. The curse: they are immune to disease, but they can carry any disease they have touched.

\noindent\textbf{Stats}: Tier I‑II, Cap 2. Corporeal, Contagious, Immune.\\
\textbf{Harm}: 3. \textbf{Fatigue}: Body 2.

\noindent\textbf{Traits}: 
\emph{Contagion Aura} – all living creatures within Near range suffer −1 die to Endurance rolls. \emph{Infectious Touch} – touch a target; they test Endurance (DV 4) or contract a disease of your choice. \emph{Fever Dream} – nightmares prevent Fatigue recovery.

\noindent\textbf{Resolution}: Fire purges the contagion. A ritual bath in salt water blessed by a witness can suppress the aura for one day. The Well of Ichor offers a cure, but the price is a memory – always the memory of the person the Plague Bearer tried to save.

\noindent\textbf{Clock}: \emph{Plague Clock [6]} – ticks each time the Plague Bearer uses Infectious Touch or remains in a populated area for more than a day. At 3, disease spreads uncontrollably. At 6, everyone flees.

\grey{The Plague Bearer is not a monster. She is a woman who held her daughter’s hand while the girl died of fever. The fever is still in her hand. She cannot let go.}

\subsection*{The Changeling (Faceless)}

\noindent\textbf{Description}: A person who wanted to escape their past – a criminal, a disgraced noble, a broken lover. The curse: they can change appearance at will, but they forget who they were originally.

\noindent\textbf{Stats}: Tier II, Cap 3. Corporeal, Shapechanger, Unstable.\\
\textbf{Harm}: 3. \textbf{Fatigue}: Body 2.

\noindent\textbf{Traits}: 
\emph{Change Face} – assume the appearance of any humanoid you have seen. Gain +2 dice to Deception. \emph{Identity Leak} – once per session, GM may have you forget a minor detail. You cannot recall it until you feed on a memory (or steal a face).

\noindent\textbf{Resolution}: A witnessed statement of the Changeling’s original name. If someone who knew them before the curse speaks their true name aloud in the presence of a witness, the Changeling will be forced to revert to their original face for one scene. If they do this three times, the curse breaks.

\noindent\textbf{Clock}: \emph{Identity Clock [6]} – ticks each time the Changeling changes faces. At 3, forget minor personal details. At 6, forget your original face entirely.

\dusana{I knew a Changeling who had been a baker. She changed faces so often that she forgot how to make bread. She baked a loaf that tasted like ash. She wept. She did not know why.}

\subsection*{The Worn – Ghouls of Zahk’Kozahn}

\noindent\textbf{Description}: Those who bear Malachai’s curses and succumb completely – who lose themselves to hunger, addiction, or despair – are not released. They are drawn back to Zahk’Kozahn, where they become the Worn. The cursed city is not empty.

\noindent\textbf{Types}:
\begin{itemize}
    \item \textbf{Shamblers}: Slow, relentless. They feel no pain and cannot be reasoned with. They exist only to feed on warmth. Harm 3, Body 3.
    \item \textbf{Screamers}: Emaciated creatures that announce their presence with a shriek that paralyses the unwary (Resolve test DV 3 or lose next action). They hunt in packs. Harm 2, Spirit 3.
    \item \textbf{Stalkers}: Silent, patient. They follow prey for hours, days, learning patterns. When they strike, they do not miss. Harm 4, Stealth 4.
    \item \textbf{Cracklers}: Their joints snap audibly as they move. They are the most “alive” of the Worn – they can still speak, still bargain, still beg. Their desperation is a weapon. Harm 3, Presence 3.
\end{itemize}

\noindent\textbf{Resolution}: The Worn are not mindless. They can be negotiated with, manipulated, even pitied. But they cannot be saved – not by mortal means. The only release is the Well of Ichor, or final death. A silver blade through the heart will kill any Worn permanently, but the silver will turn black and cannot be used again.

\noindent\textbf{Clock}: \emph{Zahk’Kozahn Calls [8]} – ticks each time a Worn is encountered. When full, the party is drawn into the cursed city. The only exit is the Well.

\lerris{I spoke to a Crackler in the tunnels beneath the city. She remembered her name. She remembered her daughter’s face. She asked me to cut her throat. I did not. I left her there. I write this now as a confession. The name was Alayse. It was my mother’s name. The Crackler was not my mother. But she remembered.}

\section*{Curses as Architecture (GM Tool)}

\noindent A curse from Malachai is not a spell; it is a **rule inserted into the world**. Write it like a clause in a contract:

\emph{“Until a silver coin that never spends is returned to the salt pan, the debtor cannot drink fresh water. The water will turn to brine in their mouth.”}

\noindent\textbf{Give it a clock}: 
\emph{Interest [6]} – ticks each time the debtor ignores the curse. When full, the curse spreads to a loved one.

\noindent\textbf{Provide a counter}: 
Return the coin, or find a witness who saw the original bargain. The witness must speak the debtor’s name aloud at the salt pan.

\noindent\textbf{SB Menu (Curses)}:
\begin{itemize}
    \item 1 SB: minor echo (old order resurfaces – a forgotten debt is remembered).
    \item 2 SB: warped threshold (the curse affects a door, a window, a bridge).
    \item 3 SB: layer slip (the curse becomes visible to bystanders – social consequences).
    \item 4 SB: ninth propagation (a second curse appears, related to the first).
\end{itemize}

\saikou{I have seen a curse spread from a father to his son. The son did not know why his water tasted of salt. He learned. The learning was the curse.}

\section*{Quick Reference – Cursed & Corrupted}

\begin{tabular}{|l|l|l|l|}
\hline
\textbf{Creature} & \textbf{Tier} & \textbf{Vulnerability} & \textbf{Resolution} \\
\hline
Linn Sea‑Wolf & III & Silver & Kill or Well of Ichor \\
Ykrul Sky‑Hound & III & Silver, confession & Kill or confession \\
Valewood Root‑Walker & III & Fulfil hospitality & Kill or fulfil promise \\
Aeler Stone‑Hide & III & Silver, keystone clause & Kill or correct violation \\
Lesser Vampire & III & Wooden stake, fire, sun & Kill or Well of Ichor \\
Slasher & II‑III & Silver, fire, first victim’s name & Kill or speak the name \\
Memory Eater & II & False memories & Starve or feed falsehood \\
Plague Bearer & I‑II & Fire, salt bath & Kill or suppress \\
Changeling & II & True name spoken thrice & Break curse \\
Worn (any) & I‑III & Silver (single use) & Kill or Well of Ichor \\
\hline
\end{tabular}

\lerris{I keep this table on the inside cover of my ledger. The ledger has nine pages. The ninth page is for the curse I have not yet named. When I name it, I will write it here. The ink will be warm.}

\grey{Malachai does not curse. She offers. The curse is the acceptance.}
\saikou{The cure is real. The freedom is an illusion. The Well is a waiting room.}
\dusana{I have seen a man drink from the Well and forget his daughter’s name. He said it was worth it. I did not believe him. I do not believe him now.}

